% Section B.3: the answers to the parts of lectures/L03/appendix/c_exercises.md that are not code.
% Every testbench message quoted was produced by a design built to Chapters 1 to 5 and broken in
% exactly the way the part describes.

\section{Chapter 3: The Synchronizer and the Receiver}
\label{sol:c3}
The a) parts of both exercises are code: \code{sync}, checked by \code{sync_tb}, and
\code{uart_rx}, checked by \code{uart_rx_tb}.

\solution{c3:ex:1}{\code{sync}}
\ans{b} The first two-edge check fails and the run stops:

\begin{console}
sync: output changed after only one clock edge!
\end{console}

\noindent It says exactly what was deleted: the input reached the output one edge after it changed,
so there is one register between them where there should be two.

Every other simulation in the course would pass, because \textbf{nothing in a simulator is ever
metastable}. A one-flop synchronizer is, functionally, a correct synchronizer with one cycle less
latency, and no other testbench cares about one cycle: \code{uart_rx_tb} drives \code{rx_s2}
directly and never instantiates \code{sync}, and in \code{uart_top_tb} the looped-back line changes
in step with the same clock, so it never violates anything. The bug is purely electrical: on the
bench, the single flop samples a genuinely asynchronous line, sometimes inside its setup and hold
window, and its output occasionally hangs between levels for long enough to reach the receiver's
logic. Only a testbench that checks the latency itself can see the missing flop, which is why
\code{sync_tb} checks it.

\ans{c} Cases 0 to 4 still pass: the reset is held low across two clock edges there, so a synchronous
reset has time to clear the chain. Case 5 asserts \code{reset_s2_n} a quarter of a period after an
edge with \code{1011} standing on the wide output and looks one nanosecond later:

\begin{console}
sync: reset_s2_n did not clear the output asynchronously!
\end{console}

\noindent The assertion has to be asynchronous so that the chain clears the moment reset is asserted,
whether or not a clock edge follows: a design is reset when its clock has not started yet, or has
stopped, and every other block in \file{hw/} clears the same instant, so a block that waited for an
edge would be the one block out of step. The release, on the other hand, is not \code{sync}'s
problem. \code{reset_s2_n} comes from \code{reset_sync}, which already lines its rising edge up with
the clock, two flip-flops after \code{reset_n} rises. By the time the release reaches \code{sync}
it is an ordinary synchronous signal, so testing it in the sensitivity list is safe.

\ans{d} A module that needs a signal as an input cannot also be where that signal comes from.
\code{sync} resets its flops from \code{reset_s2_n}; to produce \code{reset_s2_n} it would have to
be reset by its own output, which is undefined until the module has already run, and the raw
\code{reset_n} would have to go through its data input instead, where the assertion would take two
clock edges to arrive and would not arrive at all without a clock. \code{reset_sync} is built the
other way round: the raw \code{reset_n} is the asynchronous clear of both its flops, and a constant
\code{'1'} is clocked through them, so the assertion needs no clock and the release is synchronous.
The two modules look alike, two flops in series, and are wired in opposite directions, so the course
keeps both.

\ans{e} Each bit of a wide \code{sync} is its own two-flop chain. If two input bits change on the same
edge of their own domain and the change lands inside the first flops' setup and hold window, each
first flop may go metastable and each resolves independently, one to the new value and one to the
old. The one that resolved to the new value appears at the output two edges later, the other three
edges later, and for one clock \code{sync_out} shows one bit's new value beside the other's old one.

That is not a bug in \code{sync}: it promises that every bit it delivers is a clean level, and that
each input change arrives within two or three edges. It does not promise that bits which changed
together arrive together, and nothing built from independent flops can. For eight unrelated pins it
does not matter, since each pin is read on its own. For a counter it is fatal: going from
\code{0111} to \code{1000}, every bit changes, and a mixed sample can be any of sixteen values,
including \code{1111} and \code{0000}, which the counter never held. A multi-bit value crosses
domains as a Gray code, where one bit changes per step, or through a handshake or an asynchronous
FIFO, never through a wide synchronizer.

\solution{c3:ex:2}{\code{uart_rx}}
\ans{b} \textbf{The stop bit} drifts furthest. The receiver resynchronizes only on the start edge; every
later sample is timed by counting its own ticks from there, so a rate error accumulates with the
time since that edge. The first data bit is sampled about 1.6 bit periods after the edge, the stop
bit about 9.7, so the stop bit carries six times the drift.

With the transmitter faster by a fraction $f$, its bit period is $1/(1+f)$ of the receiver's. Its
stop bit, the tenth bit of the frame, ends ten of its own bit periods, 160 of its ticks, after the
edge. The stop-bit sample has to land before that.

\begin{itemize}
  \item \textbf{An ideal tick-8 sampler} samples the stop bit at the centre of its tenth bit, 9.5 bit
    periods, 152 of its ticks, after the edge. It needs $152 < 160/(1+f)$, so $f < 160/152 - 1 =
    5.3\,\%$. Counting the up to one tick a 16-times sampler inherently takes to notice the edge,
    $f < 160/153 - 1 = 4.6\,\%$.
  \item \textbf{This receiver} samples the stop bit 155 ticks after the tick that noticed the edge
    (9 to the start re-check, 17 to the first data bit, 16 for each of the other seven, 17 to the
    stop bit), and the edge is noticed up to one tick after it happens, so up to 156 ticks after the
    edge. It needs $156 < 160/(1+f)$, so $f < 160/156 - 1 = 2.6\,\%$.
\end{itemize}

\noindent So the transmitter can run about 5\,\% fast against the ideal sampler, and only about
2.5\,\% fast against this one, which spends half the budget on its own late sampling phase before
any mismatch starts. The budget is shared by both ends: the DE0-CV's divider alone already runs
0.47\,\% fast at 115200. The same late phase buys margin the other way: a transmitter up to
$155/144 - 1 = 7.6\,\%$ \emph{slower} still has its stop bit under the sample.

\ans{c} With the re-check removed, a single tick of low on an idle line starts a frame. The receiver
then samples a line that has already gone back to idle, reads eight \code{1}s and a high stop bit,
and delivers a perfectly well-formed \code{0xFF} with \code{valid} and no framing error. Simulated
with a one-tick glitch, the receiver without the re-check delivered one byte, \code{0xFF}, and no
framing error; with the re-check it delivered nothing. Note that \code{uart_rx_tb} still passes
without the re-check: none of its three cases contains a glitch, so the bench cannot see the
difference.

It is more than theoretical because real lines are noisy, and the moments they are noisiest are the
ones nobody is watching: a cable plugged in or pulled out, the far end powering up or down, a long
unshielded run beside a motor or a switching supply, a receive pin left floating with no pull-up.
Each of those can put a short low pulse on an idle line, and without the re-check each pulse is a
phantom \code{0xFF} in the RX FIFO.

\ans{d} \code{0x53} is \code{0101_0011}, so least significant first the bits arrive in this order:

\begin{narrowtable}{@{}lcccccccc@{}}
  \thead{Arrival} & 1st & 2nd & 3rd & 4th & 5th & 6th & 7th & 8th \\
  \midrule
  \thead{Value} & \code{1} & \code{1} & \code{0} & \code{0} & \code{1} & \code{0} & \code{1} &
    \code{0} \\
  \thead{\code{bit_idx}, index stored} & 0 & 1 & 2 & 3 & 4 & 5 & 6 & 7 \\
\end{narrowtable}

\noindent so \code{frame(7 downto 0)} holds \code{0}, \code{1}, \code{0}, \code{1}, \code{0},
\code{0}, \code{1}, \code{1} from bit 7 down, which is \code{0101_0011}, \code{0x53}. Counting down
from 7 instead stores the first bit in bit 7, so \code{frame} holds \code{1100_1010}: \code{data_out}
reads \code{0xCA}, the \textbf{bit reversal} of the byte sent, and \code{uart_rx_tb} says so:

\begin{console}
uart_rx_tb: received byte mismatch, expected 0x53, got 0xCA!
\end{console}
